% Options for packages loaded elsewhere
\PassOptionsToPackage{unicode}{hyperref}
\PassOptionsToPackage{hyphens}{url}
%
\documentclass[
  11pt,
]{article}
\usepackage{amsmath,amssymb}
\usepackage{iftex}
\ifPDFTeX
  \usepackage[T1]{fontenc}
  \usepackage[utf8]{inputenc}
  \usepackage{textcomp} % provide euro and other symbols
\else % if luatex or xetex
  \usepackage{unicode-math} % this also loads fontspec
  \defaultfontfeatures{Scale=MatchLowercase}
  \defaultfontfeatures[\rmfamily]{Ligatures=TeX,Scale=1}
\fi
\usepackage{lmodern}
\ifPDFTeX\else
  % xetex/luatex font selection
\fi
% Use upquote if available, for straight quotes in verbatim environments
\IfFileExists{upquote.sty}{\usepackage{upquote}}{}
\IfFileExists{microtype.sty}{% use microtype if available
  \usepackage[]{microtype}
  \UseMicrotypeSet[protrusion]{basicmath} % disable protrusion for tt fonts
}{}
\makeatletter
\@ifundefined{KOMAClassName}{% if non-KOMA class
  \IfFileExists{parskip.sty}{%
    \usepackage{parskip}
  }{% else
    \setlength{\parindent}{0pt}
    \setlength{\parskip}{6pt plus 2pt minus 1pt}}
}{% if KOMA class
  \KOMAoptions{parskip=half}}
\makeatother
\usepackage{xcolor}
\usepackage[letterpaper,margin=2.5cm]{geometry}
\usepackage{graphicx}
\makeatletter
\def\maxwidth{\ifdim\Gin@nat@width>\linewidth\linewidth\else\Gin@nat@width\fi}
\def\maxheight{\ifdim\Gin@nat@height>\textheight\textheight\else\Gin@nat@height\fi}
\makeatother
% Scale images if necessary, so that they will not overflow the page
% margins by default, and it is still possible to overwrite the defaults
% using explicit options in \includegraphics[width, height, ...]{}
\setkeys{Gin}{width=\maxwidth,height=\maxheight,keepaspectratio}
% Set default figure placement to htbp
\makeatletter
\def\fps@figure{htbp}
\makeatother
\setlength{\emergencystretch}{3em} % prevent overfull lines
\providecommand{\tightlist}{%
  \setlength{\itemsep}{0pt}\setlength{\parskip}{0pt}}
\setcounter{secnumdepth}{-\maxdimen} % remove section numbering
\ifLuaTeX
\usepackage[bidi=basic]{babel}
\else
\usepackage[bidi=default]{babel}
\fi
\babelprovide[main,import]{spanish}
% get rid of language-specific shorthands (see #6817):
\let\LanguageShortHands\languageshorthands
\def\languageshorthands#1{}
% Ajustes de calidad editorial para los boletines IES.
% Se incluyen con: pandoc --include-in-header=assets/editorial.tex
% Nota: pandoc ya carga `microtype` automáticamente si está disponible,
% por eso aquí solo se afinan los pies de figura.
\usepackage{caption}     % control de los pies de figura
\captionsetup{font=small, labelfont=bf}
\ifLuaTeX
  \usepackage{selnolig}  % disable illegal ligatures
\fi
\IfFileExists{bookmark.sty}{\usepackage{bookmark}}{\usepackage{hyperref}}
\IfFileExists{xurl.sty}{\usepackage{xurl}}{} % add URL line breaks if available
\urlstyle{same}
\hypersetup{
  pdflang={es},
  hidelinks,
  pdfcreator={LaTeX via pandoc}}

\author{}
\date{}

\begin{document}

\hypertarget{termodinuxe1mica-financiera-boletuxedn-toc-insights}{%
\section{TERMODINÁMICA FINANCIERA: BOLETÍN TOC
INSIGHTS}\label{termodinuxe1mica-financiera-boletuxedn-toc-insights}}

\textbf{Análisis Sistémico de PEMEX: La Propuesta de Escisión bajo la
Teoría de Restricciones}

\href{https://es-us.noticias.yahoo.com/proponen-escindir-pemex-negocios-deficitarios-050356273.html}{Escisión
de PEMEX}

\begin{center}\rule{0.5\linewidth}{0.5pt}\end{center}

\hypertarget{titular-y-resumen-ejecutivo}{%
\subsubsection{1. Titular y Resumen
Ejecutivo}\label{titular-y-resumen-ejecutivo}}

La propuesta analizada por \emph{El Economista} de escindir los negocios
deficitarios de Petróleos Mexicanos (PEMEX) evalúa a la empresa bajo los
dogmas de la contabilidad de costos corporativa privada. Desde la
perspectiva de la Teoría de Restricciones (TOC), desmembrar la
infraestructura petrolera ignorando su función sistémica como
amortiguador macroeconómico destruye el \textbf{Throughput} real del
país, confundiendo la optimización contable aislada con la maximización
de la soberanía y el bienestar social.

\begin{center}\rule{0.5\linewidth}{0.5pt}\end{center}

\hypertarget{diagnuxf3stico-del-sistema-seguxfan-toc}{%
\subsubsection{2. Diagnóstico del Sistema según
TOC}\label{diagnuxf3stico-del-sistema-seguxfan-toc}}

\begin{itemize}
\item
  \textbf{Restricción del Sistema (Cuello de Botella):} El verdadero
  cuello de botella no radica en la existencia de negocios secundarios,
  sino en una restricción dual: la insuficiencia de inversión focalizada
  en la exploración rentable combinada con un flujo de caja estrangulado
  por compromisos fiscales e ineficiencias logísticas.
\item
  \textbf{Parámetros Clave del Sistema:}
\item
  \textbf{Throughput (\(T\)):} El valor energético y estabilidad
  macroeconómica provistos al aparato productivo nacional para blindarlo
  frente a choques externos (como crisis geopolíticas).
\item
  \textbf{Inversión / Inventario (\(I\)):} Activos físicos (refinerías,
  pozos, red logística) y la deuda financiera consolidada.
\item
  \textbf{Gasto de Operación (\(OE\)):} Costos operativos, nómina y
  mantenimiento necesarios para sostener el andamiaje corporativo y
  operativo.
\item
  \textbf{Nube de Evaporación del Conflicto:}
\end{itemize}

\begin{figure}
\centering
\includegraphics{Diagramas/nube-evaporacion-pemex.pdf}
\caption{Nube de Evaporación del Conflicto}
\end{figure}

\begin{center}\rule{0.5\linewidth}{0.5pt}\end{center}

\hypertarget{dicotomuxeda-clave-1-contabilidad-de-costos-tradicional-vs.-contabilidad-del-throughput-ta}{%
\subsubsection{3. Dicotomía Clave 1: Contabilidad de Costos Tradicional
vs.~Contabilidad del Throughput
(TA)}\label{dicotomuxeda-clave-1-contabilidad-de-costos-tradicional-vs.-contabilidad-del-throughput-ta}}

\begin{itemize}
\tightlist
\item
  \textbf{Visión de la Contabilidad de Costos Tradicional:} Analiza las
  refinerías o los complejos petroquímicos de forma fragmentada. Si una
  unidad presenta pérdidas contables o márgenes negativos, exige su
  escisión, venta o cierre para ``mejorar el EBITDA'' y reducir costos
  por unidad.
\item
  \textbf{Visión de la Contabilidad del Throughput (TA):} Reconoce que
  recortar componentes operativos locales bajo el pretexto de eficiencia
  aislada puede colapsar el flujo global de valor. Un activo con
  aparente déficit contable puede funcionar como el amortiguador
  indispensable que protege el \emph{Throughput} sistémico de toda la
  economía nacional frente a la volatilidad de precios internacionales.
\end{itemize}

\begin{center}\rule{0.5\linewidth}{0.5pt}\end{center}

\hypertarget{dicotomuxeda-clave-2-teoruxeda-monetaria-moderna-mmt-vs.-teoruxeda-de-restricciones-toc}{%
\subsubsection{4. Dicotomía Clave 2: Teoría Monetaria Moderna (MMT)
vs.~Teoría de Restricciones
(TOC)}\label{dicotomuxeda-clave-2-teoruxeda-monetaria-moderna-mmt-vs.-teoruxeda-de-restricciones-toc}}

\begin{itemize}
\tightlist
\item
  \textbf{Visión MMT:} Al tratarse de un Estado con soberanía monetaria,
  las restricciones financieras nominales son secundarias; los apoyos
  fiscales o déficits en PEMEX no están limitados por ingresos
  tributarios previos, sino por la capacidad real de recursos y el
  riesgo inflacionario resultante.
\item
  \textbf{Visión TOC:} El dinero es meramente un medio de facilitación;
  la restricción fundamental es física y operativa. Incluso con
  soberanía monetaria, si la capacidad de extracción y refinación (la
  restricción real del sistema energético) está subfinanciada o mal
  coordinada por desvíos de recursos hacia el pago de deuda o gasto
  corriente en lugar de pozos productivos, la creación de dinero o
  transferencias fiscales solo incrementan el inventario financiero
  (\(I\)) y generan ruido inflacionario sin acelerar el flujo físico de
  bienestar.
\end{itemize}

\begin{center}\rule{0.5\linewidth}{0.5pt}\end{center}

\hypertarget{conclusiuxf3n-y-prescripciuxf3n-estratuxe9gica-5-pasos-de-enfoque-de-goldratt}{%
\subsubsection{5. Conclusión y Prescripción Estratégica (5 Pasos de
Enfoque de
Goldratt)}\label{conclusiuxf3n-y-prescripciuxf3n-estratuxe9gica-5-pasos-de-enfoque-de-goldratt}}

\begin{enumerate}
\def\labelenumi{\arabic{enumi}.}
\tightlist
\item
  \textbf{Identificar la Restricción:} Reconocer que el cuello de
  botella no es la amplitud de la empresa pública, sino la distorsión en
  la asignación de recursos (dinero destinado a sanear pasivos o deuda
  en lugar de elevar la capacidad de los campos rentables).
\item
  \textbf{Explotar la Restricción:} Maximizar la eficiencia operativa de
  los activos actuales de extracción y refinación para asegurar el
  suministro interno frente a la volatilidad geopolítica.
\item
  \textbf{Subordinar el Resto del Sistema:} Alinear los esquemas
  fiscales y las transferencias gubernamentales para que no estrangulen
  la liquidez operativa indispensable de la empresa.
\item
  \textbf{Elevar la Restricción:} Inyectar capital de inversión (\(I\))
  de manera directa en la perforación y modernización tecnológica de los
  campos de alta productividad, ampliando la capacidad del amortiguador
  energético.
\item
  \textbf{Prevenir la Inercia:} Evitar la aplicación ciega de dogmas
  corporativos privados que dictan la fragmentación de empresas
  estatales sin evaluar el impacto sistémico sobre la estabilidad
  macroeconómica y el bienestar social.
\end{enumerate}

\end{document}
